\section{The initial run}\label{sec:initial}

In this section the initial run has an arbitrary law $(\nu_\plus,\nu_\minus)$ that does not
depend on $N$, with $D<\infty$ almost surely. For $0\le k\le N$ let
\[
 \mathcal D^\plus_N(k)=N\!\!\sum_{k/2\le d\le\min(k,N-1)}\!\!\nu_\plus(d)\,F^\minus_{N-d}(k-d),
\]
\[
 \mathcal D^\minus_N(k)=N\!\!\sum_{(N-k)/2\le d\le\min(N-k,N-1)}\!\!\nu_\minus(d)\,F^\plus_{N-d}(k).
\]
The term $\mathcal D^\plus_N(k)$ is the event that the initial run is a plus run that ends
after at least half of the $k$ plus steps. The rest of the walk is then a fresh minus walk
that has few plus steps, so it sits in its own end layer, where the probabilities are of
order $T(N)$ and not $1/N$.

\begin{introtheorem}[the initial run]\label{thm:initial}
Assume (BAL), (SPAN) and (SRT$_\pm$).
\begin{enumerate}[(a)]
\item (The defect.) Uniformly over $k$ with $k/N$ in each compact subset of $(0,1)$,
\[
 NP_N(k)=\La(k/N)+\mathcal D^\plus_N(k)+\mathcal D^\minus_N(k)+o(1).
\]
\item (A necessary condition.) If $NP_N(\lfloor x_0N\rfloor)\to\La(x_0)$ for one
$x_0\in(0,1)$, then $nT_\plus(n)\,\nu_\pm(n)\to0$ as $n\to\infty$.
\item (A sufficient condition.) Assume moreover (LB$_\pm$). If $n\nu_\pm(n)\to0$, then
$NP_N(k)=\La(k/N)+o(1)$ uniformly over $k$ with $k/N$ in each compact subset of
$(0,1)$. This holds in particular when $\nu_\plus$ and $\nu_\minus$ are eventually
nonincreasing.
\item (An exact criterion.) Assume moreover (LB$_\pm$) and $u_\pm(i)>0$ for all $i\ge0$.
Then the conclusion of (c) holds if and only if, for all $0<x_1<x_2<1$,
\begin{align*}
 &\sup_{x_1N\le k\le x_2N}NT_\plus(N)\sum_{0\le i\le k/2}\nu_\plus(k-i)\,\omega_\plus(i)\to0
 \qquad\text{and}\\
 &\sup_{x_1N\le k\le x_2N}NT_\plus(N)\sum_{0\le i\le(N-k)/2}\nu_\minus(N-k-i)\,\omega_\minus(i)\to0 .
\end{align*}
\item (A finite initial run that changes the bulk.) Let $c>0$, $n_i=2^{2^i}$, and let
$\nu_\plus(n_i)=c/(n_iT_\plus(n_i))$ for $i\ge i_0$, with the rest of the mass of the
initial run placed anywhere. Then for every $x_0\in(0,1)$, along the $N$ with
$\lfloor x_0N\rfloor\in\{n_i\}$,
\[
 \liminf NP_N(\lfloor x_0N\rfloor)\ge\La(x_0)+c\,b\,x_0^{\alpha-1}(1-x_0)^{-\alpha}.
\]
Under the hypotheses of (d) the condition of (b) is not sufficient. An initial run with
a cluster of $\lfloor\sqrt{n_i}\rfloor$ equal masses $T_\plus(\sqrt{n_i})/(n_iT_\plus(n_i))$
just below each $n_i$, for $i\ge i_0$, and the rest of its mass at $d=1$, satisfies the
condition of (b), while the local law fails along the
$N$ with $\lfloor x_0N\rfloor\in\{n_i\}$.
\item (The end layer.) Assume only (BAL) and (SPAN). For each fixed $j\ge0$,
\[
 P_N(N-j)=Q_\plus(N)\ones_{\{j=0\}}+T_\plus(N)\kappa_\plus(j)
 +\sum_{1\le m<N/2}\nu_\plus(N-m)\,\Pi^\minus_m(j)+o(T_\plus(N)),
\]
where
\begin{equation}\label{eq:kappa-prep}
 \kappa_\plus(j)=\pi_\plus u_\minus(j)\ones_{\{j\ge1\}}+\sum_{d=1}^{j}\nu_\minus(d)\,u_\minus(j-d).
\end{equation}
Hence $P_N(N-j)/T_\plus(N)\to\kappa_\plus(j)$ for every $j\ge1$ if and only if
\[
 W^\plus_N:=\sum_{1\le m<N/2}\nu_\plus(N-m)\,T_\plus(m)=o(T_\plus(N)).
\]
This holds if $n\nu_\plus(n)\to0$, and it holds for every fresh start.
\end{enumerate}
The same statements hold with the signs exchanged.
\end{introtheorem}

Part (b) is the first place where the infinite mean shows. For $1<\alpha<2$ the bulk
needs no condition on the initial run on any macroscopic interior~\cite{paperF}.
Here a first run of length close to $k$ leaves a fresh walk with very few plus steps,
whose probability is of order $T(N)\gg1/N$. In (f) the window condition replaces the
condition $\nu_\plus=o(T_\plus)$ of the finite-mean case, because now
$\sum_{m<N/2}T_\plus(m)$ grows like $NT_\plus(N)$. At $j=0$ the end value is
$P_N(N)=Q_\plus(N)$ exactly. So in (f) the profile value at $j=0$ is
$\lim_NQ_\plus(N)/T_\plus(N)$ when this limit exists, and not $\kappa_\plus(0)=0$.

We compare the conditions. The necessary condition (b) asks that $\nu_\pm(n)$ be
$o(1/(nT_\plus(n)))$, and the sufficient condition (c) asks that $\nu_\pm(n)$ be $o(1/n)$.
The gap between them is the factor $T_\plus(n)$, which tends to $0$. The exact criterion (d)
says that what matters is a weighted window sum of $\nu_\pm$ just below the point, with the
weight $\omega_\pm(i)\asymp1/(iT(i))$, whose partial sums up to $w$ grow like $1/T(w)$. A single
spike of size $c/(nT(n))$ at $n$ fails (b). A cluster of $w$ equal masses $T(w)/(nT(n))$ just
below $n$ satisfies (b) and still fails (d), because the weights $\omega(i)$ over the cluster add
up to about $1/(\alpha T(w))$. Under (LB$_\pm$) every initial
run with a regularly varying tail $\Prob(D\ge n)$ of negative index and an eventually monotone
mass function satisfies (c), whatever its index, so it changes neither the bulk nor, by (f),
the end layers. Note that the initial run may be much heavier than the other runs. For instance
$\nu_\pm(n)\propto n^{-1-\gamma}$ with $0<\gamma<\alpha$ gives an initial run with a heavier tail
than every later run, and under (LB$_\pm$) the local Lamperti law still holds. The effect of such a run is
only felt at speed $\Prob(D\ge\varepsilon N)$, as the computations of Section~\ref{sec:verify}
show for $\gamma=0.2$ and $\alpha=1/2$.

\subsection{Proofs}

\begin{proof}[Proof of (a)]
Fix $0<x_1<x_2<1$ and let $k\in[x_1N,x_2N]$. In Proposition~\ref{prop:first-run} the atoms
vanish. Consider the plus part $\sum_{d<N}\nu_\plus(d)F^\minus_{N-d}(k-d)$, where only
$d\le k$ contributes. Fix a small $\varepsilon>0$. For $d<\varepsilon N$ the horizon is
$n=N-d\ge(1-\varepsilon)N$ and $(k-d)/n\in[x_1/2,x_2]$. Theorem~\ref{thm:local} for the
fresh minus walk holds uniformly in the horizon $n\to\infty$ and in the ratio in a
compact set. Since $|(k-d)/n-k/N|\le2\varepsilon$ and $\La$ is uniformly continuous on
$[x_1/2,x_2]$, this gives $NF^\minus_{N-d}(k-d)=\La(k/N)+O(\varepsilon)+o(1)$. So the
terms with $d<\varepsilon N$ contribute $\Prob(\Sigma=\plus,D<\varepsilon N)\La(k/N)+O(\varepsilon)+o(1)
=\pi_\plus\La(k/N)+O(\varepsilon)+o(1)$. For $\varepsilon N\le d<k/2$ the ratio
$(k-d)/(N-d)$ lies in $[x_1/2,x_2]$ and $N-d\ge N/2$, so $NF^\minus_{N-d}(k-d)\le C$ by
Theorem~\ref{thm:local}, and these terms contribute at most $C\Prob(D\ge\varepsilon N)\to0$.
The terms with $k/2\le d\le k$ form $\mathcal D^\plus_N(k)$. The minus part
$\sum_{d<N}\nu_\minus(d)F^\plus_{N-d}(k)$ is the same, with $d\le N-k$ and the ratio
$k/(N-d)\le2x_2/(1+x_2)<1$ for $d<(N-k)/2$. Since $\pi_\plus+\pi_\minus=1$, letting
$N\to\infty$ and then $\varepsilon\to0$ proves (a).
\end{proof}

\begin{proof}[Proof of (b)]
All terms of $\mathcal D^\pm_N$ are nonnegative. Since $F^\minus_{N-k}(0)=T_\minus(N-k)$, the
term $d=k$ of $\mathcal D^\plus_N(k)$ is $N\nu_\plus(k)T_\minus(N-k)$. In the same way the
term $d=N-k$ of $\mathcal D^\minus_N(k)$ is $N\nu_\minus(N-k)T_\plus(k)$. By (a) and the hypothesis both tend
to $0$ along $k=k_N=\lfloor x_0N\rfloor$. Both $k_N$ and $N-k_N$ take every large integer
value, and $NT_\minus(N-k_N)$ and $NT_\plus(k_N)$ are comparable with $k_NT_\plus(k_N)$ and
$(N-k_N)T_\plus(N-k_N)$ by (R2). Hence $nT_\plus(n)\nu_\pm(n)\to0$.
\end{proof}

\begin{proof}[Proof of (c) and (d)]
Write $i=k-d$, so $\mathcal D^\plus_N(k)=N\sum_{0\le i\le k/2}\nu_\plus(k-i)F^\minus_{N-k+i}(i)$.
The horizon $n=N-k+i$ lies between $(1-x_2)N$ and $N$, and
$i/n\le(k/2)/(N-k/2)\le x_2/(2-x_2)<1$. Theorem~\ref{thm:ends}(b) at the minus end, with
the signs exchanged, applies to the fresh minus walk under (LB$_\minus$) and (SRT$_\plus$).
It gives, uniformly in these $i$,
\[
 F^\minus_n(i)=T_\minus(n)\bigl[u_\plus(i)\Phi_{1/b}(i/n)+o(\omega_\plus(i))\bigr].
\]
Since $u_\plus\le C\omega_\plus$, $\Phi_{1/b}$ is bounded on $[0,x_2/(2-x_2)]$ and
$T_\minus(n)\le CT_\plus(N)$, we get
$\mathcal D^\plus_N(k)\le C\,NT_\plus(N)\sum_{0\le i\le k/2}\nu_\plus(k-i)\omega_\plus(i)$. If
moreover $u_\plus(i)>0$ for all $i$, then $u_\plus\ge c\,\omega_\plus$ by (SRT$_\plus$), and
$\Phi_{1/b}$ is bounded below on that interval, so the reverse inequality holds with a
smaller constant for large $N$. The same holds for $\mathcal D^\minus_N$. Since
$NP_N(k)-\La(k/N)=\mathcal D^\plus_N(k)+\mathcal D^\minus_N(k)+o(1)$ uniformly with
$\mathcal D^\pm_N\ge0$, this proves (d).

For (c), let $\nu_\plus(d)\le\epsilon/d$ for $d\ge x_1N/2$. Then by (R3)
\[
 NT_\plus(N)\sum_{0\le i\le k/2}\nu_\plus(k-i)\omega_\plus(i)\le\frac{2\epsilon NT_\plus(N)}k
 \Bigl(1+\sum_{1\le i\le k}\frac1{iT_\plus(i)}\Bigr)\le C\epsilon\,\frac{NT_\plus(N)}{kT_\plus(k)}\le C_{x_1}\epsilon .
\]
The same bound holds for the minus sum, so $\mathcal D^\pm_N\to0$ uniformly and (a) gives
the local law. If $\nu_\plus$ is eventually nonincreasing, then
$\frac n2\nu_\plus(n)\le\sum_{n/2<d\le n}\nu_\plus(d)\to0$.
\end{proof}

\begin{proof}[Proof of (e)]
The masses $c/(n_iT_\plus(n_i))$ are summable, since $1/(nT_\plus(n))$ is regularly varying
with the negative index $\alpha-1$ and $n_i$ grows doubly exponentially. So the initial run
is well defined for $i_0$ large. Let $k=\lfloor x_0N\rfloor=n_i$. By (a) and the
nonnegativity of the other terms,
$NP_N(k)\ge\La(k/N)+N\nu_\plus(k)T_\minus(N-k)+o(1)$. Now
\[
 N\nu_\plus(k)T_\minus(N-k)=c\,\frac Nk\,\frac{T_\minus(N-k)}{T_\plus(k)}\to
 c\,x_0^{-1}\,b\Bigl(\frac{1-x_0}{x_0}\Bigr)^{-\alpha}=c\,b\,x_0^{\alpha-1}(1-x_0)^{-\alpha},
\]
by (R1). For the second example put $n=n_i$ and $w=\lfloor\sqrt n\rfloor$, and let
$\nu_\plus(d)=T_\plus(w)/(nT_\plus(n))$ for $n-w<d\le n$. The mass of the cluster is
$wT_\plus(w)/(nT_\plus(n))$, which is regularly varying in $n$ with index $-(1-\alpha)/2<0$,
so the clusters have summable mass. For $d$ in the cluster,
$dT_\plus(d)\nu_\plus(d)\le CT_\plus(w)\to0$, so the condition of (b) holds. At $k=n$ the
first sum in (d) is at least
\[
 NT_\plus(N)\,\frac{T_\plus(w)}{nT_\plus(n)}\sum_{0\le i<w}\omega_\plus(i)
 \ge c\,\frac{NT_\plus(N)}{nT_\plus(n)}\cdot T_\plus(w)\cdot\frac1{\alpha T_\plus(w)}(1+o(1))\ge c'>0
\]
by (R3), so the criterion of (d) fails and the local law fails along these $N$.
\end{proof}

\begin{proof}[Proof of (f)]
By Proposition~\ref{prop:first-run}, with $F^\minus_{N-d}(N-j-d)=\Pi^\minus_{N-d}(j)$,
\[
 P_N(N-j)=Q_\plus(N)\ones_{\{j=0\}}+\sum_{d<N}\nu_\plus(d)\,\Pi^\minus_{N-d}(j)
 +\sum_{1\le d\le j}\nu_\minus(d)\,F^\plus_{N-d}(N-j),
\]
since $F^\plus_{N-d}(N-j)=0$ for $d>j$. By Theorem~\ref{thm:ends}(a) for a fresh plus walk,
the last sum is $T_\plus(N)\sum_{d\le j}\nu_\minus(d)u_\minus(j-d)+o(T_\plus(N))$. In the
middle sum, split at $d=N/2$. For $d\le N/2$, Theorem~\ref{thm:ends}(a) for a fresh minus
walk gives $\Pi^\minus_n(j)=T_\plus(n)[u_\minus(j)\ones_{\{j\ge1\}}+o(1)]$ uniformly in
$n\ge N/2$. Since $1\le T_\plus(N-d)/T_\plus(N)\le C$ for $d\le N/2$ and the ratio tends to
$1$ for each fixed $d$, dominated convergence gives
$\sum_{d\le N/2}\nu_\plus(d)T_\plus(N-d)=\pi_\plus T_\plus(N)(1+o(1))$. So the terms with
$d\le N/2$ give $T_\plus(N)\pi_\plus u_\minus(j)\ones_{\{j\ge1\}}+o(T_\plus(N))$. The terms
with $d>N/2$ are exactly $\sum_{1\le m<N/2}\nu_\plus(N-m)\Pi^\minus_m(j)$ with $m=N-d$. This
proves the formula.

For the criterion, let $j\ge1$. By Theorem~\ref{thm:ends}(a) for a fresh minus walk,
$\Pi^\minus_m(j)\le C_jT_\plus(m)$ for all $m$, so $W^\plus_N=o(T_\plus(N))$ makes the sum
$o(T_\plus(N))$. Conversely, suppose that $P_N(N-j)/T_\plus(N)\to\kappa_\plus(j)$ for every
$j\ge1$. Then the sum, which is nonnegative, is $o(T_\plus(N))$ for every $j\ge1$. Pick
$j^*\ge1$ with $u_\minus(j^*)>0$. Then $\Pi^\minus_m(j^*)\ge cT_\plus(m)$ for $m\ge m^*$, so
$\sum_{m^*\le m<N/2}\nu_\plus(N-m)T_\plus(m)=o(T_\plus(N))$. For each $m<m^*$ take $j=m$.
Then $\Pi^\minus_m(m)=T_\minus(m)>0$, so $\nu_\plus(N-m)=o(T_\plus(N))$. Together these give
$W^\plus_N=o(T_\plus(N))$.

If $n\nu_\plus(n)\to0$, then $\nu_\plus(N-m)=o(1/N)$ uniformly in $m<N/2$, and
$W^\plus_N\le o(1/N)\sum_{m<N/2}T_\plus(m)=o(T_\plus(N))$ by (R3). For a fresh minus start
$\nu_\plus=0$. For a fresh plus start, fix $M$. The terms $m<M$ give at most
$M\max_{n\ge N-M}f_\plus(n)=o(T_\plus(N))$, and the terms $M\le m<N/2$ give at most
$T_\plus(M)\Prob(N/2<X_\plus\le N-M)\le T_\plus(M)T_\plus(N/2)\le CT_\plus(M)T_\plus(N)$. Let
$N\to\infty$ and then $M\to\infty$.
\end{proof}

If the law of the initial run depends on $N$ and puts mass of order $1$ on lengths of
order $N$, the local law itself changes. A first run of length $zN$ is followed by a fresh
walk of length $(1-z)N$, and the local law becomes a mixture of rescaled Lamperti densities.
Section~\ref{sec:aging} proves this, under (SRT$_\pm$) and with no condition of the type
(LB), and applies it to the process observed from a late time on.
